\documentclass[10pt,letterpaper]{article}

\usepackage[margin=0.5in]{geometry}
\usepackage[T1]{fontenc}
\usepackage{lmodern}
\usepackage{xcolor}
\usepackage{enumitem}
\usepackage{tabularx}
\usepackage[hidelinks]{hyperref}
\usepackage{titlesec}
\usepackage{microtype}

\definecolor{accent}{HTML}{145A78}
\setlength{\parindent}{0pt}
\setlength{\parskip}{0pt}
\setlist[itemize]{leftmargin=1.15em,itemsep=0.8pt,topsep=1pt,parsep=0pt}
\titleformat{\section}{\large\bfseries\color{accent}}{}{0pt}{}[\vspace{-2pt}\titlerule]
\titlespacing*{\section}{0pt}{5pt}{2pt}
\pagestyle{empty}

\newcommand{\role}[4]{%
  \textbf{#1}\hfill #2\\
  \textit{#3}\hfill \textit{#4}\vspace{-3pt}
}

\begin{document}
\fontsize{9}{10.5}\selectfont

\begin{center}
  {\Large\bfseries Piter Z. Garcia Bautista}\\[1pt]
  {\large K--12 STEM Subject Matter Expert and Inclusive CS Educator}\\[2pt]
  Rochester, NY \enspace|\enspace Available for remote part-time, freelance, and contract work\\
  \href{mailto:pzg8794@rit.edu}{pzg8794@rit.edu}
  \enspace|\enspace +1 (646) 288-3300
  \enspace|\enspace \href{https://linkedin.com/in/garciapiterz}{linkedin.com/in/garciapiterz}\\
  \href{https://github.com/pzg8794}{github.com/pzg8794}
  \enspace|\enspace \href{https://garciapiterz.com}{garciapiterz.com}
\end{center}

\section*{Summary}
Graduate educator and applied AI professional with experience designing
standards-aligned K--12 instruction in computer science, robotics, digital
fluency, mathematics, and responsible AI. Combines inclusive teaching,
curriculum development, technical writing, and graduate research in data
science and machine learning. Comfortable developing, reviewing, and revising
learner-centered content in distributed, technology-enabled teams.

\section*{Subject and Curriculum Strengths}
\begin{tabularx}{\textwidth}{@{}>{\bfseries}l X@{}}
Subject areas & Computer science, mathematics, data science, AI, robotics, digital literacy, and interdisciplinary STEM\\
Curriculum & Standards alignment, lesson and activity design, assessment development, content review, scaffolding, feedback cycles\\
Inclusive practice & Universal Design for Learning, disability inclusion, accessible communication, equitable participation\\
Technical tools & Python, SQL, R, Java, C++, JavaScript, Git/GitHub, Jupyter, LaTeX, learning and collaboration platforms\\
AI and data & Machine learning, model evaluation, data preparation, reproducible experimentation, responsible and fairness-aware AI\\
\end{tabularx}

\section*{Education and Teaching Experience}
\role{Computer Science Teacher Candidate}{2025--Present}
{University of Rochester / Pine Brook Elementary School}{Rochester, NY}
\begin{itemize}
  \item Co-design and facilitate standards-aligned learning in coding, robotics,
    digital fluency, mathematics, and responsible AI for varied K--12 learners.
  \item Develop lessons, activities, prompts, and assessments using Universal
    Design for Learning supports and revise materials from learner and mentor
    feedback.
  \item Translate technical ideas into clear explanations, examples, and
    hands-on experiences that support multiple ways of engaging and showing
    understanding.
\end{itemize}

\role{Graduate Assistant, Quantum and AI Research}{Fall 2025--Present}
{Rochester Institute of Technology, Software Engineering}{Rochester, NY}
\begin{itemize}
  \item Build and document Python experiments for adaptive quantum-network
    routing, online learning, and resource allocation.
  \item Review technical results for accuracy, create reproducible
    comparison-ready artifacts, and communicate complex methods through
    reports, presentations, and manuscript-ready materials.
\end{itemize}

\section*{Applied Technical Experience}
\role{Data Solutions Engineer}{Sep 2022--Aug 2024}
{VEDADATA Inc.}{New York, NY}
\begin{itemize}
  \item Designed maintainable Python data workflows, validation checks, and
    documentation for analytics and business applications.
  \item Collaborated across technical and nontechnical teams to clarify
    requirements, review outputs, and improve reliability and usability.
\end{itemize}

\role{AI Data Solutions Engineer}{Jan 2021--Sep 2022}
{VIOME Inc.}{New York, NY}
\begin{itemize}
  \item Developed machine-learning workflows for healthcare use cases,
    including preprocessing, feature engineering, supervised-model
    experimentation, evaluation, and clear reporting.
\end{itemize}

\section*{Education}
\textbf{M.S., Teaching Computer Science K--12}, University of Rochester
\hfill Expected December 2026\\
GPA: 3.9+; NSF Robert Noyce Scholar; Advanced Certificate in Leadership in
Disability and Inclusive Practices

\textbf{M.S., Data Science}, Rochester Institute of Technology
\hfill Expected December 2026\\
GPA: 3.9+; applied AI, bioinformatics, neural networks, and reproducible data systems

\textbf{M.S., Computer Science}, Rochester Institute of Technology
\hfill 2015\\
Concentrations in artificial intelligence, big data analytics, and computer graphics

\textbf{B.S., Computer Engineering Technology}, Farmingdale State College
\hfill 2012

\section*{Recognition}
NSF Robert Noyce Scholar (2024--2026)
\enspace|\enspace IEEE Alumni Member (2009--Present)
\enspace|\enspace RIT/MESCYT merit-based graduate funding

\end{document}
